\section{Q2: Future Representation and Granularity}
\label{app:q2}

\subsection{Representation Scope and Comparison Axes}

Q2 compares what the future representation describes and at what granularity.
The reference history encoding and information flow are provisional selections
from the preceding phases: \placeholder{MISSING: Q2 reference configuration}.
All representations use declared physical prediction times and the same
available input information within a contrast. Changes in resolution, target
count, prediction objective, and online computation are recorded, rather than
being hidden in a representation name.

The working space includes target-only, context-only, uniform-granularity, and
heterogeneous-granularity futures. ``Target'' means the manipulated object or
task-relevant local state; ``context'' includes non-target movers, obstacles,
support/receptacle geometry, and free-space constraints relevant to feasible
action. The exact region construction and deployable target assignment are
\Pending. A privileged region mask is supervision or an oracle condition, not
an unreported visual input.

\begin{table}[t]
  \centering
  \caption{Q2 comparison specification. The actual admitted representation,
  resolution, and measured costs remain pending for each row.}
  \label{tab:app_q2_scope}
  \small
  \begin{tabularx}{\columnwidth}{@{}lXX@{}}
    \toprule
    Axis & Controlled question & Required declaration \\
    \midrule
    Target-only & Is target future sufficient for this task? & Target state/region; omitted context \\
    Context-only & Can surrounding motion/constraints guide action? & Context state; current target access \\
    Uniform & What follows from one common prediction granularity? & Space/time resolution and token budget \\
    Heterogeneous & Does allocating detail differently help? & Target/context allocation and total cost \\
    Representation & Which future state exposes useful action information? & SE(3), flow, points, trajectory, RGB/latent as admitted \\
    \bottomrule
  \end{tabularx}
\end{table}

``Target fine, relevant context coarse'' is a falsifiable hypothesis. It is not
encoded as the selected architecture before the comparison. If it lacks a
stable benefit, the study reports the tasks, horizons, and budgets where the
effect is absent or negative, and chooses the final instance accordingly.
The completed candidate set and reasons for excluding any candidate are
\placeholder{MISSING: Q2 candidate admission record}.

\subsection{Horizon, Region, and Budget Matching}

Prediction horizon is measured from capture time to the time the action takes
effect, as in Eq.~\eqref{eq:app_physical_horizon}. Spatial density and temporal
density are independent axes: more points at one instant are not more future
time samples. For every candidate, the reporting record includes the physical
time grid, region support, number of predicted elements, feature dimensions,
trainable parameters, prediction NFE if applicable, and end-to-end latency.

Where exact parameter and latency equality cannot both hold, the comparison
states which budget is matched and supplies the measured alternative cost.
An advantage caused by denser temporal samples, more parameters, a longer
history, or lower action age cannot be assigned to the semantic representation
alone. Resolution/cost curves, if measured, report the actual points; missing
points are not connected into an invented curve. The accepted budget envelope
and completed resolution settings are \Pending.

The geometric candidate of Appendix~\ref{app:implementation} uses rigid pose
for the target and metric displacement for distributed motion. Point and
trajectory candidates must specify their correspondence and sampling rules;
RGB and latent candidates must specify decoding and the action-visible
interface. A common action interface can be used when that is the controlled
question, but its bottleneck capacity and training exposure must be reported.
Changing both the representation and the backbone is reserved for an explicit
interaction check rather than attributed to Q2 alone.

\subsection{Prediction Measures and Action-Level Measures}

Representation-specific errors describe whether the chosen target is predicted
well. For a rigid target, translation error and symmetry-aware rotation error
are reported separately. Flow uses valid-point EPE in metric units and can
include target rigidity consistency. Point, RGB, and latent errors are not
numerically pooled with pose error. The available per-horizon values are
\PoseADE\ cm, \RotADE$^\circ$, \FlowEPE\ cm, and
\FlowSEConsistency\ for the historical geometric fields; the full arrays and
the applicable candidate mapping remain \Pending.

The action-level comparison uses the same task outcome definitions for all
representations. Contact timing and relative contact pose connect prediction
to interception and grasping. Feasible approach, receptacle containment, and
collision/safety outcomes test context-dependent constraints. No scalar
prediction metric substitutes for task success. A lower average error could
still miss the few times or regions that determine contact.

\begin{table}[t]
  \centering
  \caption{Pending Q2 action-level expansion. Each final row is a measured
  candidate under the stated reference; this table currently records no ranking.}
  \label{tab:app_q2_results}
  \small
  \begin{tabularx}{\columnwidth}{@{}Xccc@{}}
    \toprule
    Future scope & SR (\%) & Contact error & Unsafe (\%) \\
    \midrule
    Target-only & \PendingCI & \Pending & \Pending \\
    Context-only & \PendingCI & \Pending & \Pending \\
    Uniform granularity & \PendingCI & \Pending & \Pending \\
    Heterogeneous granularity & \PendingCI & \Pending & \Pending \\
    Strong history baseline & \PendingCI & \Pending & \Pending \\
    Declared future oracle & \PendingCI & \Pending & \Pending \\
    \bottomrule
  \end{tabularx}
\end{table}

The contact-error field in Table~\ref{tab:app_q2_results} is expanded into
separate timing, translation, and orientation columns when actual results are
available; the units and aggregation must then be explicit. Task/seed/horizon
breakdowns, confidence intervals, and safety denominators are
\placeholder{MISSING: Q2 complete result arrays}. A main-paper claim about
context requires evidence on the relevant constraint, not only a lower global
flow error.

\subsection{Oracles, Calibration, and Failure Strata}

Future oracles estimate headroom under explicitly privileged information.
A future-XYZ oracle and a future-SE(3) oracle are different controls; both
remain distinct from a vision model. The historical gap-closed diagnostic is
\begin{equation}
 G_{\rm oracle}=
 \frac{\operatorname{SR}(P)-\operatorname{SR}(B)}
 {\operatorname{SR}(O)-\operatorname{SR}(B)},
 \label{eq:app_oracle_gap}
\end{equation}
for predictive method $P$, matched history baseline $B$, and declared oracle
$O$. It is interpretable only when the denominator supports a meaningful
comparison; the raw differences and intervals are reported regardless.
The pending historical field is \OracleGapClosed\%; oracle definition,
coverage, and the measured denominator remain \Pending.

Calibration is reported only for candidates with a probabilistic output and a
defined calibration procedure. Error and uncertainty can be stratified by
visibility, speed, acceleration, orientation change, horizon, and event phase
when those labels are available. A comparison before and after robot contact
also states whether it is action-conditioned. Failures that reverse a
representation ranking stay in the main argument, with the full breakdown
here: \placeholder{MISSING: Q2 negative results and conditional effects}.

After Q3 changes the interaction backbone, the leading and contrasting Q2
candidates are rechecked. Representation selection and backbone selection are
therefore provisional until the measured interactions are known, rather than
three independent choices assumed to compose.
